\clearpage
\markboth{Abstract}{Abstract}
\phantomsection
\addcontentsline{toc}{section}{Abstract}

% 标题：Abstract
\begin{center}
    \setstretch{1.5}
    \zihao{3}\bfseries\selectfont
    \vspace*{12pt}
    Abstract
    \vspace*{6pt}
\end{center}

% 摘要正文
\begingroup
    \zihao{-4}\selectfont
    \setstretch{1.25}
    \setlength{\parindent}{2em}
    \setlength{\parskip}{0pt}
    
    The geological-geochemical behaviors and seven genetic types of associated elements in coals of the differences of associated elements in coals from North and Southwest China. The distributing models and mechanism of rare earth elements in coals were discussed in detail. In this paper, the background values and five sources of PGEs in coals are present, and the distribution of PGEs is characterized by Pt-Pd pattern. By using the technology of TOF-SIMS, the study shows that the bacteria play an important role in the activation and enrichment of Cu and Zn during the formation studied in detail.
    \par
    \vspace{\baselineskip}
    \centerline{...... (Continue with approximately 1500 words) ......}
    \par
    \vspace{\baselineskip}
    The author groping investigated the effects of the synsedimentary volcanic-ash on the enrichment and the occurrence of associated elements in coal; the non-sulfide Fe and Cu were found in coal, and the petrographic names and types of the special material composed by the volcanic ash, organic matter and detrial material of terrigenous origin were schemed out. There is a high moving contrast of U, Mo, and Cu in the reductive and the oxidized barrier. The oxidized barrier leads to the high activation of above elements, and the reductive barrier leads to the secondary accumulation.
    \par
\endgroup

\begingroup
    \setstretch{1.25}
    \zihao{-4}
    \textbf{Key Words: } coal; associated elements; geological-geochemical behaviours; enrichment models
\endgroup
\clearpage
